% !TeX root = ../../../main.tex
\subsection{Discussion and Conclusion}
\label{ch08:sec:discussion}


Our results suggest that the primary benefit of multimodal alignment is temporal calibration rather than event discovery.
\ChEightRebuttal{
Across models, structured evidence changed event recovery little, while mean paired case-level temporal differences generally favored multimodal calibration: 7 of 12 model--metric comparisons across concordance and AULTC were positive and none were negative. However, uncertainty in the magnitude of these improvements was substantial given the 40-case sample; paired confidence intervals excluded zero only for the DeepSeek V3.2 AULTC improvement. Source-stratified estimates also differed between i2b2-derived and MIMIC-IV cases (\hyperref[ch08:supp:paired-bootstrap]{Supplementary Material \S\ref*{ch08:supp:paired-bootstrap}} and \hyperref[ch08:supp:dataset-stratified]{\S\ref*{ch08:supp:dataset-stratified}}).
Within the text-primary formulation, narrative text defines the candidate event set and structured rows provide partial temporal evidence; the results therefore establish the value of narrative content for events represented in the note, not that narrative is the primary source of all clinically meaningful events.}

% Adding retrieved structured EHR evidence has little effect on event match rate, but more consistently improves temporal localization. This indicates that structured rows are most useful for sharpening when text-derived events occurred, not for replacing narrative event extraction. In this sense, structured EHR data act as partial temporal evidence for narrative-derived patient trajectories, while the narrative remains the primary source of clinically meaningful event content.

The structured-data gap analysis explains why this distinction matters.
Structured EHR data is not merely a noisier version of the narrative record; it offers a different view of the patient trajectory.
When relevant counterparts exist, their timestamps can anchor timeline calibration.
However, many clinically meaningful narrative events are absent, delayed, or semantically compressed in tabular form, including symptoms, functional status, temporal qualifiers, patient-reported information, and clinical context. 
\ChEightRebuttal{These findings characterize
asymmetric coverage within the present text-primary analysis; they do not assess structured-only events or establish that either source is complete.}
% Models built only on structured data may therefore miss important patient states even when their timestamps are precise, while models built only on text may recover those states but place them less accurately in time.

The ablation results also support the scaffolded design.
One-shot timeline generation requires a model to identify clinical events, infer relative timings, resolve vague temporal expressions, and decide how to use structured evidence in a single step.
By first constructing an anchor-event scaffold, the workflow separates the temporal backbone of the encounter from supporting narrative details.
This mirrors clinical reasoning, where events such as admission, diagnostic testing, treatment initiation, deterioration, procedures, and discharge provide temporal reference points for interpreting the rest of the patient course.
Staged calibration then allows structured evidence to stabilize the backbone before refining the full trajectory.


\noindent\textbf{Limitations:} 
\ChEightRebuttal{This study has several limitations. First, event match rates remain modest, so concordance and AULTC measure temporal quality conditional on successful event recovery rather than accuracy over the complete trajectory. Matched events were clinically heterogeneous, but unmatched events more often contained later-course and granular numerical documentation (\hyperref[ch08:supp:matched-unmatched]{Supplementary Material \S\ref*{ch08:supp:matched-unmatched}}, \hyperref[ch08:tab:supp-matched-unmatched]{Table~\ref*{ch08:tab:supp-matched-unmatched}}). Second, the 40-case sample provides only high-uncertainty evidence about stability across data sources and clinical
settings; moreover, each case was annotated once by one of three medically trained annotators, without duplicate annotation or inter-annotator agreement. Third, the framework emits one continuous timestamp per event; when timing is vague or absent, it infers a point estimate from context, and the model-generated certainty flag is uncalibrated (\hyperref[ch08:supp:temporal-certainty]{Supplementary Material \S\ref*{ch08:supp:temporal-certainty}}, \hyperref[ch08:tab:supp-certainty]{Table~\ref*{ch08:tab:supp-certainty}}). Fourth, the text-primary design may omit clinically meaningful structured-only events. Structured timestamps also have source-specific semantics and may represent collection, ordering, administration, result availability, device recording, or
documentation rather than clinical onset. Finally, the seven-call workflow has substantially higher inference cost than the single-step baseline.}


% This study has several limitations to consider. \textbf{First}, event match rates remain modest, reflecting the difficulty of extracting standalone clinical events from dense discharge summaries and aligning them to manual timelines; temporal metrics should therefore be interpreted jointly with event recovery, since they are computed only on matched events. \textbf{Second}, although the evaluation includes 40 manually annotated discharge summaries, the benchmark remains modest in scale because multimodal annotation is costly: annotators must inspect and cross-reference narrative summaries and structured EHR tables. Ongoing annotation efforts aim to expand evaluation to a larger and more diverse set of cases. \textbf{Third}, the framework is text-primary: narrative events define the candidate trajectory, while structured rows mainly provide temporal evidence, so clinically meaningful structured-only events absent from the note may be missed. \textbf{Fourth}, structured timestamps require careful interpretation, since chart times, lab draw times, medication administration times, and documentation times may differ from true event onset. \textbf{Finally}, the workflow depends on instruction-tuned LLMs and structured prompts; although we use open-weight models and local inference, model choice, prompting, and runtime costs remain practical considerations.


\noindent\textbf{Practical use and clinical impact:} This framework is not intended for autonomous clinical decision-making, but for temporally grounded research and clinician-reviewed workflows. Reconstructed trajectories could support retrospective cohort construction, treatment-window definition, temporal phenotyping, leakage auditing, and downstream trajectory modeling.
\ChEightRebuttal{Future work should evaluate larger, independently annotated cohorts; support interval-valued timestamps, qualitative relations, and explicit time-unknown outputs; develop bidirectional event discovery that preserves source provenance and timestamp semantics; and test stage-specific cost-quality tradeoffs. 
Overall, the findings support a retrieval-augmented view in which narrative text supplies candidate-event semantics, structured rows provide partial and source-dependent temporal evidence, and preserve-by-default calibration can improve temporal grounding for some events when relevant evidence is available.}

% Future work should expand evaluation across institutions and diseases, assess clinician review of reconstructed timelines, and test whether temporally grounded trajectories improve downstream prediction, causal modeling, and trial-screening tasks. Overall, our findings support a retrieval-augmented view of clinical timeline reconstruction: narrative text provides the clinical event semantics, structured EHR rows provide partial temporal evidence, and aligning them conservatively yields patient trajectories that are more temporally grounded than trajectories reconstructed from either text or tables alone.
%%%%%%%%%%%%%%%%%%%%%%%%%%%%%%%%%%%%%%%%%%%%%%%%%%

